\section{作用量不变为什么给出守恒流}

取实标量场 \(\phi\) 的一阶 Lagrange 密度
\(\mathcal L(\phi,\partial_\mu\phi,x)\)，定义
\(S[\phi]=\int_\Omega\mathcal L\,\dd^dx\)。
这里 \(\Omega\) 是有分段光滑边界的时空区域。先让变分
\(\eta=\delta\phi\) 在边界为零，才可以从驻值条件得到场方程。

\begin{definition}[一阶变分与 Euler--Lagrange 表达式]
一阶变分是
\(\delta S[\phi;\eta]
=\left.\frac{d}{d\varepsilon}S[\phi+\varepsilon\eta]\right|_{\varepsilon=0}\)。
令 \(\pi^\mu=\partial\mathcal L/\partial(\partial_\mu\phi)\)，
\(E(\phi)=\partial\mathcal L/\partial\phi-\partial_\mu\pi^\mu\)。
若 \(\delta S=0\) 对所有紧支集 \(\eta\) 成立，则场满足 \(E(\phi)=0\)。
\end{definition}

链式法则与分部积分给出未丢边界项的恒等式
\[
\delta\mathcal L
=\frac{\partial\mathcal L}{\partial\phi}\eta
+\pi^\mu\partial_\mu\eta
=E(\phi)\eta+\partial_\mu(\pi^\mu\eta).
\]
对区域积分，边界项为
\(\int_{\partial\Omega}\pi^\mu\eta n_\mu\,d\Sigma\)。
紧支集变分使它消失；由变分法基本引理才得 \(E(\phi)=0\)。
若允许非零边界变分，必须另加边界条件或边界作用量。

\begin{theorem}[Noether 第一变分定理]
设一参数场变换在固定坐标 \(x\) 处的无穷小变化为
\(\eta(x)\)，且对任意场配置有
\(\delta\mathcal L=\partial_\mu B^\mu\)。
则流 \(j^\mu=\pi^\mu\eta-B^\mu\) 满足恒等式
\[
\partial_\mu j^\mu=-E(\phi)\eta.
\]
特别地，若 \(\phi\) 满足 Euler--Lagrange 方程，则
\(\partial_\mu j^\mu=0\)。
\end{theorem}
\begin{proof}
将定理假设 \(\partial_\mu B^\mu=\delta\mathcal L\) 代入刚才的分部积分恒等式，移项得
\(\partial_\mu(\pi^\mu\eta-B^\mu)=-E(\phi)\eta\)。
在场方程成立时右侧为零。这同时说明守恒是对满足运动方程的场的结论。
\end{proof}

平移给一个可复算的例子。若 \(\mathcal L\) 无显含 \(x\)，沿第 \(\nu\) 个坐标方向作主动变换，
\(\eta=-a^\nu\partial_\nu\phi\)，
\(\delta\mathcal L=-a^\nu\partial_\nu\mathcal L
=\partial_\mu(-a^\mu\mathcal L)\)。
因此
\(j^\mu=-a^\nu T^\mu{}_\nu\)，其中
\[
T^\mu{}_\nu=\pi^\mu\partial_\nu\phi-\delta^\mu{}_\nu\mathcal L
\]
是正则能动张量。由于常数 \(a^\nu\) 任意，场方程给
\(\partial_\mu T^\mu{}_\nu=0\)。
若 \(\mathcal L\) 显含 \(x^\nu\)，平移不是对称，此结论不能照搬。

Lorentz 对称也可以在同一套公式中做，而不必再猜一个守恒量。对实标量场，取上一节的无穷小坐标变换
\(\xi^\nu(x)=\omega^\nu{}_{\rho}x^\rho\)，其中
\(\omega_{\nu\rho}=-\omega_{\rho\nu}\)。主动变换在固定 \(x\) 处给
\(\eta=-\xi^\nu\partial_\nu\phi\)。若密度只由 Lorentz 标量
\(\phi\) 和 \(\partial_\mu\phi\partial^\mu\phi\) 组成，则
\(\delta\mathcal L=-\partial_\mu(\xi^\mu\mathcal L)\)：这里
\(\partial_\mu\xi^\mu=\omega^\mu{}_{\mu}=0\)，反对称性保证缩并的导数项相消。Noether 公式于是给
\(j^\mu=-\xi^\nu T^\mu{}_{\nu}\)。把任意的反对称参数提出，可取角动量流
\[
M^\mu{}_{\rho\sigma}
=x_\rho T^\mu{}_{\sigma}-x_\sigma T^\mu{}_{\rho}.
\]
例如 \(\mathcal L=\tfrac12\partial_\mu\phi\partial^\mu\phi-\tfrac12\kappa^2\phi^2\) 时，
\(T_{\rho\sigma}=\partial_\rho\phi\partial_\sigma\phi-\eta_{\rho\sigma}\mathcal L\) 对称。直接求散度得到
\[
\partial_\mu M^\mu{}_{\rho\sigma}
=T_{\rho\sigma}-T_{\sigma\rho}
+x_\rho\partial_\mu T^\mu{}_{\sigma}
-x_\sigma\partial_\mu T^\mu{}_{\rho}=0
\]
（最后一步使用场方程）。一般带自旋的场，其正则能动张量不必对称，角动量流还须加由场表示产生的自旋项；所以这里的纯轨道表达式只对所讨论的标量场直接成立。

内部相位变换给另一种守恒流，而且可以把每一项算出。设
\(\phi\) 为复标量，取常数 \(\kappa>0\)（量纲为长度的倒数）及
\[
\mathcal L=\partial_\mu\phi^*\partial^\mu\phi-\kappa^2\phi^*\phi .
\]
在变分时把 \(\phi,\phi^*\) 当作两个独立实变量的等价记法。
常数参数 \(\varepsilon\) 的变换
\(\delta\phi=\ii\varepsilon\phi\)、
\(\delta\phi^*=-\ii\varepsilon\phi^*\)
使 \(\mathcal L\) 严格不变，故 \(B^\mu=0\)。
两个场的动量分别是
\(\pi_\phi^\mu=\partial^\mu\phi^*\)、
\(\pi_{\phi^*}^\mu=\partial^\mu\phi\)。
从 Noether 公式得到每单位 \(\varepsilon\) 的流
\[
j^\mu=\ii\bigl(\phi\,\partial^\mu\phi^*
-\phi^*\,\partial^\mu\phi\bigr).
\]
运动方程是 \((\Box+\kappa^2)\phi=0\) 及其复共轭。
直接求散度时，两项一阶导数的乘积互相抵消，余下
\[
\partial_\mu j^\mu
=\ii(\phi\Box\phi^*-\phi^*\Box\phi)=0.
\]
这一步比仅说“相位对称所以电荷守恒”多明确了场方程在何处使用。
若让 \(\varepsilon\) 暂时依赖 \(x\)，则
\(\delta\mathcal L=j^\mu\partial_\mu\varepsilon\)；
由任意紧支集 \(\varepsilon(x)\) 的变分再一次得到
\(\partial_\mu j^\mu=0\)。局域化参数是检验守恒流的计算技巧，
不意味着原作用量具有局域规范对称。

局部守恒还不自动等于整个区域的荷守恒。沿用上一节的 \(x^0=ct\)，令
\(Q(t)=\int_\Sigma j^0(ct,\vb{x})\,\dd^{d-1}x\)。由
\(\partial_0j^0+\partial_i j^i=0\) 和 \(\partial_t=c\partial_0\)，散度定理给
\(\dd Q/\dd t=-\int_{\partial\Sigma}\vb{J}\cdot\dd\vb{S}\)，其中物理空间通量 \(\vb{J}=c(j^1,\ldots,j^{d-1})\)。
只有边界通量为零、周期边界相消，或无穷远处衰减足够快，\(Q\) 才不随时间变化。

\begin{exercise}
取 \(\mathcal L=\tfrac12\partial_\mu\phi\,\partial^\mu\phi-\tfrac12\kappa^2\phi^2\)，其中 \(\kappa>0\) 为常数。
求 \(E(\phi)\) 与 \(T^\mu{}_\nu\)，并直接计算
\(\partial_\mu T^\mu{}_\nu=-E(\phi)\partial_\nu\phi\)。
\end{exercise}
